\par\bigskip
\noindent\textbf{6. Singularidades cuadráticas no degeneradas}

\medskip
El objeto de esta sección es estudiar las deformaciones de un
$k$-esquema cuyas únicas singularidades sean singularidades
cuadráticas no degeneradas. En vista de 4.13, esto se convierte, en
una situación adecuada, en un problema puramente local.

Volvemos a la notación de \S4. Así, $\Lambda$ es un anillo local
noetheriano completo fijo con cuerpo residual $k$, y
$\underline C_\Lambda$ denota la categoría de las $\Lambda$-álgebras
locales artinianas con cuerpo residual $k$.

Sea $X$ un $k$-esquema de tipo finito y sea $x$ un punto cerrado de
$X$. Recordemos que $x$ se denomina una \emph{singularidad cuadrática
no degenerada racional sobre $k$} si
\[
\widehat{\mathcal O}_{X,x}
 \simeq k[[x_1,\ldots,x_n]]/(f)
\]
y
\[
\left(\frac{\partial f}{\partial x_1},\ldots,
      \frac{\partial f}{\partial x_n}\right)
=(x_1,\ldots,x_n).
\]
Siempre se puede tomar $f$ como una forma cuadrática. En efecto:

\medskip
\noindent\textbf{Lema 6.1.}
Sea $f\in k[[x_1,\ldots,x_n]]$, y pongamos
$M=(x_1,\ldots,x_n)$. Si
\[
\left(\frac{\partial f}{\partial x_1},\ldots,
      \frac{\partial f}{\partial x_n}\right)=M,
\]
entonces existen elementos $x_1',\ldots,x_n'$ que satisfacen
$x_i'\equiv x_i\pmod{M^2}$ tales que
\[
f=q(x_1',\ldots,x_n')
\]
para una forma cuadrática $q$.

\smallskip
\noindent\textit{Demostración.}
Utilizamos aproximaciones sucesivas. Escribamos
\[
f=f_2+h_3,
\]
donde $f_2$ es cuadrática y $h_3\in M^3$. Como las derivadas parciales
de $f$ generan $M$, existen
$\varepsilon_1^{(1)},\ldots,\varepsilon_n^{(1)}\in M^2$ tales que
\[
\sum_i\varepsilon_i^{(1)}\frac{\partial f}{\partial x_i}
 \equiv h_3\pmod{M^4}.
\]
Pongamos $x_i^{(1)}=x_i-\varepsilon_i^{(1)}$. Entonces
\[
\begin{aligned}
f(x_1^{(1)},\ldots,x_n^{(1)})
&=f(x_1,\ldots,x_n)
 -\sum_i\varepsilon_i^{(1)}\frac{\partial f}{\partial x_i}
 +\cdots\\
&\equiv f_2(x_1,\ldots,x_n)\pmod{M^4}.
\end{aligned}
\]
Así,
\[
f(x_1^{(1)},\ldots,x_n^{(1)})
=f_2(x_1,\ldots,x_n)+h_4,
\qquad h_4\in M^4.
\]
Además, las derivadas respecto de las nuevas variables siguen
generando $M$. Por tanto, se pueden elegir
$\varepsilon_1^{(2)},\ldots,\varepsilon_n^{(2)}\in M^3$ tales que
\[
\sum_i\varepsilon_i^{(2)}
 \frac{\partial f}{\partial x_i^{(1)}}
 \equiv h_4\pmod{M^5}.
\]

%%%%%%%%%%  scan idx 136  |  folio 125  %%%%%%%%%%

Se sigue que
\[
\begin{aligned}
&f(x_1^{(1)}-\varepsilon_1^{(2)},\ldots,
  x_n^{(1)}-\varepsilon_n^{(2)})\\
&\qquad =f(x_1^{(1)},\ldots,x_n^{(1)})
 -\sum_i\varepsilon_i^{(2)}
       \frac{\partial f}{\partial x_i^{(1)}}+\cdots\\
&\qquad \equiv f_2(x_1,\ldots,x_n)\pmod{M^5}.
\end{aligned}
\]
Prosiguiendo de este modo, se encuentran elementos
$\varepsilon_i^{(\nu)}\in M^{\nu+1}$ tales que, poniendo
\[
x_i^{(\nu)}
=x_i-\bigl(\varepsilon_i^{(1)}+\cdots+
            \varepsilon_i^{(\nu)}\bigr),
\]
se tiene
\[
f(x_1^{(\nu)},\ldots,x_n^{(\nu)})
\equiv f_2(x_1,\ldots,x_n)\pmod{M^{\nu+3}}.
\]
Sea $\varepsilon_i=\sum_\nu\varepsilon_i^{(\nu)}$. La completitud da
entonces
\[
f(x_1-\varepsilon_1,\ldots,x_n-\varepsilon_n)
=f_2(x_1,\ldots,x_n),
\]
o, equivalentemente,
\[
f(x_1,\ldots,x_n)
=f_2(x_1+\varepsilon_1,\ldots,x_n+\varepsilon_n).
\]
Tomando $x_i'=x_i+\varepsilon_i$ se demuestra la afirmación.
\hfill$\square$
